\begin{scanpage}{0234}{216}
\noindent the current becomes noticeable also in those directions where, according to physical optics, the fringing field must be equal to zero or have a small value.

The problem of diffraction of a plane wave with its incidence on a cone along its axis (\S\ 17) serves as a clear example of how important the above-indicated factors are. Although in this case the nonuniform part of the current, concentrated near the cone's vertex, has practically no influence on the scattering, nevertheless, the physical optical approach gives values for the radar thickness which are tens of decibels smaller than the experimental values, even with large dimensions of the cone. The deciding factor here is the nonuniform part of the current flowing in the vicinity of the sharp circular base rim of the conical surface; the nonuniform part of the current has an especially large value for sharply pointed cones.

Another interesting example of a similar nature is the scattering of a plane wave by a finite paraboloid of rotation (\S\ 18) where the physical optics approach leads to qualitatively incorrect results. The effective scattering area calculated in this approach turns out to be a periodic function of the paraboloid length, and with certain lengths it becomes zero which most certainly does not correspond to reality.

The investigation of the diffraction of edge waves shows (Chapter V) that for flat plates one may limit oneself to consideration of secondary diffraction, if their linear dimensions are one-and-a-half to two times larger than the wavelength.

Let us note that we attempted to obtain equations for the scattering characteristics which would possess physical visualizability and which would be convenient for making calculations. In keeping with this, we were obliged to introduce various kinds of interpolation equations and simplified equations which satisfy the formulated requirements, but in return are not in the general case the dominant terms of the rigorous asymptotic expansion in powers of the small parameter $\lambda/a$.
\end{scanpage}
